\subsection{Dedekind 整环的定义}

设 A 为整环。显然以下条件等价：

\begin{enumerate}
\def\labelenumi{(\alph{enumi})}
\item
  A 的非零素理想中任意两个在包含关系下都不可比较；
\item
  \(\mathbf{A}\) 的非零素理想都是极大理想；
\item
  A 的非零素理想都具有高度1。
\end{enumerate}

定义1. 非零素理想全为极大理想的 Krull 整环称为 Dedekind 整环。

\textbf{Dedekind 整环的例子}\par

\begin{enumerate}
\def\labelenumi{(\arabic{enumi})}
\item
  每个主理想整环都是 Dedekind 整环。
\item
  设 K 是 Q 的有限扩张，A 是 K 中 Z 的整闭包。环 A 是 Krull 整环（第1节，第8号，命题12）。设 p 是 A 的非零素理想。则 p ∩ Z 非零（第V章，第2节，第1号，命题1的推论），因而是 Z 的极大理想；所以 p 是 A 的极大理想（同上，命题1）。因此，A 是 Dedekind 整环。一般来说，A 不是主理想整环（《代数》第VII章，第1节，习题12）。
\item
  * 设 V 是一个仿射代数簇，A 是 V 上正则函数的环。假设 A 不是域（即 V 不退化为一点）。A 为 Dedekind 整环的充要条件是 V 是没有奇点的不可约曲线：因为 A 是整环等价于 V 不可约；A 的每个非零素理想都是极大理想等价于 A 是一条曲线；最后，由于 A 是诺特环，A 是 Krull 整环等价于 A 是整闭的，也就是说 V 是正规曲线，或者说 V 没有奇点。 *
\item
  分式环 \(S^{-1}A\)（其中 A 是 Dedekind 整环）在 \(0 \notin S\) 时是 Dedekind 整环。因为 \(S^{-1}A\) 是 Krull 整环（第1节，第4号，命题6），并且根据第II章，第2节，第5号，命题11，\(S^{-1}A\) 的每个非零素理想都是极大理想。
\end{enumerate}

